\documentclass[10pt,conference]{IEEEtran}
\usepackage{amsmath}
\usepackage{amssymb}
\usepackage{graphicx}
\usepackage{booktabs}
\usepackage{hyperref}

\title{Policy Gradient Methods for HW3: REINFORCE, Baselines, and DQN Comparison}
\author{Taha Majlesi}

\begin{document}
\maketitle

\begin{abstract}
We study policy gradient methods across four tasks: a custom GridWorld, CartPole with and without a value-function baseline, MountainCarContinuous with Gaussian policies, and a CartPole comparison between REINFORCE and DQN. We summarize objectives, variance-reduction strategies, and discrete vs. continuous action handling. All experiments are configured with centralized seeds and metadata saved alongside the notebooks.
\end{abstract}

\section{Introduction}
Policy gradients directly optimize stochastic policies by ascending the expected return. This report aligns the accompanying notebooks with the assignment description and documents configurations, seeds, and metadata needed for reproducibility.

\section{Methodology}
\subsection{REINFORCE Objective}
For a trajectory $\tau = (s_0, a_0, r_1, \ldots, s_T)$ sampled from $\pi_\theta$, the objective is
\begin{equation}
J(\theta) = \mathbb{E}_{\tau \sim \pi_\theta}\left[\sum_{t=0}^{T-1} \gamma^t r_{t+1}\right],
\end{equation}
with gradient
\begin{equation}
\nabla_\theta J(\theta)=\mathbb{E}_{\tau}\left[\sum_{t=0}^{T-1} \nabla_\theta \log \pi_\theta(a_t|s_t)\, G_t\right],
\end{equation}
where $G_t$ is the return from $t$ onward.

\subsection{Baseline Variance Reduction}
Subtracting a baseline $b(s_t)$ yields
\begin{equation}
\nabla_\theta J(\theta)=\mathbb{E}_{\tau}\left[\sum_{t=0}^{T-1} \nabla_\theta \log \pi_\theta(a_t|s_t)\, (G_t-b(s_t))\right],
\end{equation}
leaving the estimator unbiased because $\mathbb{E}[\nabla_\theta \log \pi_\theta(a_t|s_t)b(s_t)]=0$.

\subsection{Gaussian Policy for Continuous Control}
For MountainCarContinuous, we use $\pi_\theta(a|s)=\mathcal{N}(\mu_\theta(s), \sigma_\theta^2)$ with log-probability
\begin{equation}
\log \pi_\theta(a|s) = -\tfrac{1}{2}\left(\frac{(a-\mu_\theta(s))^2}{\sigma_\theta^2} + 2\log \sigma_\theta + \log 2\pi\right).
\end{equation}

\subsection{DQN Baseline}
For comparison on CartPole, we implement a Deep Q-Network that minimizes the temporal-difference loss
\begin{equation}
\mathcal{L}_{\text{DQN}} = \mathbb{E}\left[\left(r + \gamma \max_{a'} Q_{\bar{\theta}}(s', a') - Q_\theta(s,a)\right)^2\right],
\end{equation}
where $\bar{\theta}$ denotes target-network parameters.

\section{Experiments}
\textbf{Environments:} GridWorld (custom), CartPole-v1 (discrete), MountainCarContinuous-v0 (continuous).\\
\textbf{Seeds/Device:} Centralized seed $42$ for Parts 1--3 and $2024$ for Part 4; device auto-detected with CUDA flag logged.\\
\textbf{Artifacts:} Each notebook writes \texttt{run\_info.json} capturing seed, device, commit hash, and hardware summary.\\
\textbf{Evaluation:} Episode returns and learning curves; policy entropy for REINFORCE variants; TD loss for DQN.

\section{Ablations}
We recommend (i) entropy regularization strength sweep, (ii) baseline on/off for CartPole, (iii) Gaussian policy std initialization for MountainCarContinuous, and (iv) replay buffer size and target update frequency for DQN.

\section{Conclusion}
Centralized configuration, deterministic seeding, and recorded metadata improve comparability across policy-gradient and value-based baselines. The accompanying notebooks implement these safeguards for reproducible HW3 results.

\appendices
\section{Hyperparameters}
\begin{table}[h]
\centering
\begin{tabular}{@{}l l@{}}
\toprule
Seed & 42 (P1--P3), 2024 (P4) \\
Discount $\gamma$ & per notebook defaults \\
Learning rate & per notebook defaults \\
Episodes & per notebook defaults \\
\bottomrule
\end{tabular}
\end{table}

\section{Reproducibility Metadata}
Each run stores \texttt{run\_info.json} with commit hash, hardware (machine, processor, platform), device, and CUDA availability.

\bibliographystyle{IEEEtran}
\bibliography{refs}

\end{document}


